\section{SmallWood Programming Model}
\label{sec:smallwood-model}
\label{sec:smallwood}

This section gives the SmallWood formulation of the ARC-SPRUCE NIZKs. The prover commits to packed witness rows, represented as low-degree row polynomials over a packing set \(\Omega\) and masked on a disjoint set \(\Omega'\). The verifier then checks PACS-style constraints on those commitments. The main text records only the proof interface and the semantic issuance and showing statements. Concrete row layouts, gadget realizations, and proof details are deferred to Appendix~\ref{app:smallwood}.

\subsection{Proof interface}
\label{subsec:smallwood-interface}
Fix a packing set \(\Omega=\{\omega_1,\ldots,\omega_s\}\subset \Fq\) and a disjoint hiding set
\(\Omega'=\{\omega'_1,\ldots,\omega'_\ell\}\subset \Fq\setminus\Omega\). A compiled witness is a matrix
\[
W\in\Fq^{n\times s}.
\]
For each row \(i\), the prover interpolates a polynomial \(P_i\in\Fq[X]\) of degree \(<s+\ell\) such that \(P_i(\omega_j)=W_{i,j}\) for all \(\omega_j\in\Omega\), while \(P_i(\omega'_h)\) is random on \(\Omega'\). The committed witness is therefore the family
\[
P=(P_1,\ldots,P_n),
\]
committed through the DECS/LVCS/PCS stack.

SmallWood proves PACS constraint families on these committed rows. Parallel constraints are slotwise identities
\[
f_a\!\bigl(P_1(\omega),\ldots,P_n(\omega),\Theta_{a,1}(\omega),\ldots,\Theta_{a,t_a}(\omega)\bigr)=0
\qquad\text{for every }\omega\in\Omega,
\]
where the \(\Theta_{a,j}\) are public helper rows. Aggregated constraints are domain-summed identities
\[
\sum_{\omega\in\Omega}
g_b\!\bigl(P_1(\omega),\ldots,P_n(\omega),\Psi_{b,1}(\omega),\ldots,\Psi_{b,u_b}(\omega)\bigr)=0.
\]
In ARC-SPRUCE, local algebra, replay bridges, and PRF checks are parallel; global carry or range checks are aggregated.

\subsection{Issuance as a SmallWood statement}
\label{subsec:smallwood-issuance}
\label{subsec:presign-compiled}
The public issuance statement is
\[
(\ComPublicParam,\vecc,T,\vecr_0^{\Issuer},\vecr_1^{\Issuer},\mA,B),
\]
and the hidden witness is
\[
(\vecmu,\veck,\vecr_0^{\Holder},\vecr_1^{\Holder},\bar r).
\]
Because the target \(T\) is public, the compiled issuance proof is the simpler one. It commits rows encoding the message, the hidden key, the holder randomness, and the centered randomness, together with auxiliary rows whose exact realization is deferred to the appendix.

Semantically, the proof establishes three claims. First, the public commitment
\(\vecc\) opens to the same hidden message-and-randomness witness used
everywhere else. Second, the centered rows \(R_0,R_1\) are derived correctly
from the holder and issuer contributions. Third, the public target \(T\)
matches that same hidden witness under the retained BB-tran relation.

Write \(M_\Sigma:=M+K\), and let \(M_{\Sigma}R_1\) and \(R_0R_1\) denote the
auxiliary product rows carried only for the BB-tran cleared identity. Writing
\(T^\Theta\), \((B_3B_1)^\Theta\), and \((B_3B_2)^\Theta\) for the public
replay lifts, and writing \(\widehat M_\Sigma\), \(\widehat R_0\),
\(\widehat R_1\), \(\widehat{M_{\Sigma}R_1}\), and \(\widehat{R_0R_1}\) for
the exact replay images of the corresponding source rows, the target claim is
checked as
\[
B_3^\Theta\Had T^\Theta
-T^\Theta\Had \widehat R_1
-(B_3B_1)^\Theta\Had\widehat M_\Sigma
-(B_3B_2)^\Theta\Had \widehat R_0
+B_1^\Theta\Had\widehat{M_{\Sigma}R_1}
+B_2^\Theta\Had\widehat{R_0R_1}
-1=0.
\]
Thus the issuance proof certifies one coherent hidden opening witness for both \(\vecc\) and the public target \(T\).

\subsection{Showing as a SmallWood statement}
\label{subsec:smallwood-showing}
\label{subsec:showing-compiled}
The public showing statement is the ARC verification context
\[
(\vecmu,\prftag,\mA,B,\BoundMsg,\BoundSig)
\]
together with the chosen nonce. The hidden witness consists of a signature witness \(\vecu\), the hidden key \(\veck\), the hidden randomness \((R_0,R_1)\), and the auxiliary data needed to express the proof in replay space.

The compiled showing proof commits three semantic layers. The first is a
cert-side signature surface proving shortness of \(\vecu\) and binding
\[
T=\mA\vecu.
\]
The second is the retained non-sign witness surface. It keeps the message and
randomness compressed through carrier rows, carries explicit source blocks for
\(T\), and adds the BB-tran auxiliary products
\[
M_{\Sigma}R_1,\ R_0R_1,
\]
where \(M_\Sigma=M+K\) is the packed signed message, public selector
constraints bind the message part to \(\vecmu\), and \(K\) carries the hidden
signed key. The third is the replay witness: for each replay block \(b\), the
proof carries exact replay-image rows
\[
\widehat T_b,\ \widehat m_b,\ \widehat R_{0,b},\ \widehat R_{1,b},\ \widehat{M_{\Sigma}R_1}_b,\ \widehat{R_0R_1}_b.
\]

Writing \(B_{i,b}^\Theta:=\mathcal F_{\mathrm{rep},b}(B_i)\) and likewise for
\((B_3B_1)_b^\Theta\) and \((B_3B_2)_b^\Theta\), the proof enforces the exact
source-to-replay bridges
\[
\widehat T_b=\mathcal F_{\mathrm{rep},b}(T),\qquad
\widehat m_b=\mathcal F_{\mathrm{rep},b}(M_\Sigma),
\]
\[
\widehat R_{0,b}=\mathcal F_{\mathrm{rep},b}(R_0),\qquad
\widehat R_{1,b}=\mathcal F_{\mathrm{rep},b}(R_1),
\]
\[
\widehat{M_{\Sigma}R_1}_b=\mathcal F_{\mathrm{rep},b}(M_{\Sigma}R_1),\qquad
\widehat{R_0R_1}_b=\mathcal F_{\mathrm{rep},b}(R_0R_1),
\]
and the blockwise replay residual
\[
B_{3,b}^\Theta\Had\widehat T_b
-\widehat T_b\Had\widehat R_{1,b}
-(B_3B_1)_b^\Theta\Had\widehat m_b
-(B_3B_2)_b^\Theta\Had\widehat R_{0,b}
+B_{1,b}^\Theta\Had\widehat{M_{\Sigma}R_1}_b
+B_{2,b}^\Theta\Had\widehat{R_0R_1}_b
-1=0.
\]
Hence the direct certified statement is a replay-space statement about one
hidden witness: the same short signature yields \(T\), the same compressed
carrier witness determines \(M_\Sigma,R_0,R_1\), the same auxiliary products
match those source values, and the same hidden key is used by the PRF relation
below.

\subsection{PRF companion relation}
\label{subsec:prf-checkpointed}
\label{subsec:prf-constraints}
The public tag is tied to the same hidden key carried by \(K\) through a
companion PRF relation. In the main text only its semantic parts matter:
selector constraints bind the PRF input lanes to the signed key, nonlinear
checkpoint constraints certify the S-box steps of the permutation, and final
feed-forward plus truncation bind the public pair \((\nonce,\prftag)\). The
exact checkpoint schedule is an instantiation detail and is deferred to
Appendix~\ref{app:prf} and Appendix~\ref{app:smallwood}.

\subsection{Hiding and statement strength}
\label{subsec:statement-strength}
The proof hides a family of row polynomials, not just a tuple of scalars. First, each row is masked by random values on \(\Omega'\), so unopened witness values are statistically hidden. Second, DECS/LVCS/PCS and the PACS batching expose only masked linear combinations and a few random-point openings outside \(\Omega\). Thus the verifier learns only the evaluations needed to check the constraints, and zero knowledge reduces to the hiding guarantees of the SmallWood commitment-opening stack and the simulator for the PACS layer.

The replay-space statement above is the only statement certified directly by the showing proof. The bridge back to the coefficient-domain relation used later is the following theorem suite.

\begin{lemma}[Exact replay transform]\label{lem:replay-transform}
The public replay map
\[
\mathcal F_{\mathrm{rep}}:\Rq\to\mathcal T
\]
is additive, multiplicative, and injective.
\end{lemma}

\begin{lemma}[Replay-space certification implies the cleared coefficient-domain relation]\label{lem:showing-to-cleared}
Assume the compiled showing relation of Section~\ref{subsec:showing-compiled}. Then
\[
B_3T-TR_1-(B_3B_1)(M+K)-(B_3B_2)R_0+B_1(M_{\Sigma}R_1)+B_2(R_0R_1)-1=0
\]
in \(\Rq\). Since the same relation also enforces \(T=\mA\vecu\), it follows that
\[
B_3\mA\vecu-(\mA\vecu)R_1-(B_3B_1)(M+K)-(B_3B_2)R_0+B_1(M_{\Sigma}R_1)+B_2(R_0R_1)-1=0
\]
in \(\Rq\).
\end{lemma}


\begin{corollary}[Abstract BB-tran consequence under admissibility]\label{cor:showing-to-abstract}
Assume the compiled showing relation of Section~\ref{subsec:showing-compiled}. Then
\[
M_{\Sigma}R_1=(M+K)R_1
\qquad\text{and}\qquad
R_0R_1=R_0R_1.
\]
If moreover \(B_3-R_1\in\Rq^\times\), then
\[
T=h^{\mathrm{bb\mbox{-}tran}}_{M+K,(R_0,R_1)}(B),
\]
so in particular
\[
\mA\vecu = h^{\mathrm{bb\mbox{-}tran}}_{M+K,(R_0,R_1)}(B).
\]
\end{corollary}

\begin{corollary}[Consequence of an accepting post-sign proof]\label{cor:showing-proof-consequence}
If the post-sign NIZK is knowledge sound for the compiled showing relation,
then every accepting proof yields a witness
\((\vecu,T,M,K,R_0,R_1,M_{\Sigma}R_1,R_0R_1)\) satisfying the cleared
coefficient-domain relation above and, under admissibility, the abstract
BB-tran target relation.
\end{corollary}
